\parttitle{IV}{Implementation Handbook: Cirq, Qiskit, and PennyLane}

\section{The code that ships with this document}
\label{sec:code-overview}

Every number and figure in this document is produced by the programs in the
\code{code/} directory, summarised in Table~\ref{tab:code-files}. The
environment is Python~3.12 with NumPy~2.4, Cirq~1.7, Qiskit~2.5 with
Qiskit-Aer~0.17, and PennyLane~0.45; each script runs standalone in seconds to
minutes on one CPU core. The pedagogical order of this part is deliberate:
first an \emph{exact} density-matrix implementation that defines ground truth
and hosts the validation suite; then three framework implementations, each
chosen to teach the thing that framework does best --- Cirq for transparent
circuit construction and exact expectation values, Qiskit for everything
between an algorithm and a physical device (transpilation, noise, shots,
dynamic circuits), and PennyLane for batched evaluation and differentiability
at the reservoir/variational boundary.

\begin{table*}[h]
\centering\small
\caption{Programs shipped with this document. All paths relative to
\code{code/}.}
\label{tab:code-files}
\begin{tabular}{@{}llp{8.2cm}@{}}
\toprule
File & Role & Contents \\
\midrule
\code{qrc\_core.py} & reference & exact Fujii--Nakajima reservoir, ESN baseline, ridge, NARMA, memory capacity, shot-noise model, validation suite \\
\code{datasets.py} & data & solar and load surrogates (labelled synthetic), real ENSO series \\
\code{eda.py} & analysis & all Part~VI exploratory figures and statistics \\
\code{experiments.py} & results & temporal-edge sweep, memory capacity, prediction demos, shot-noise study (Part~VII figures) \\
\code{cirq\_qrc\_minimal.py} & framework & restart-protocol reservoir in Cirq (Listing~\ref{lst:cirq}) \\
\code{qiskit\_qrc\_hw.py} & framework & hardware-aware pipeline and mid-circuit reset demo (Listing~\ref{lst:qiskit}) \\
\code{pennylane\_qrc\_hybrid.py} & framework & broadcast feature map and differentiable encoding tuning (Listing~\ref{lst:pennylane}) \\
\code{figstyle.py} & support & publication figure style \\
\bottomrule
\end{tabular}
\end{table*}

\section{An exact reference implementation}
\label{sec:reference-impl}

The reference implementation represents $\rho$ explicitly as a
$2^n\times 2^n$ complex matrix and applies Eq.~\eqref{eq:fn-update} literally.
Its purpose is threefold: it defines the analytic ceiling against which every
shot-noise and hardware-noise result is measured; it is the only
implementation in which conceptual errors cannot hide behind framework
abstractions; and it hosts the validation suite that every later change must
pass. The injection map is eleven lines
(Listing~\ref{lst:core-inject}); note that the input qubit is fixed as the
\emph{leftmost} tensor factor, so re-insertion by a Kronecker product lands in
the correct slot without any index permutation --- the single most common bug
in hand-rolled QRC code is silently rebuilding the state with the qubit order
scrambled.

\begin{strip}
\lstinputlisting[caption={The Fujii--Nakajima injection,
Eq.~\eqref{eq:fn-update}, in \code{qrc\_core.py}. \code{trace\_out\_first}
implements $\tr_1$ by reshaping to a rank-4 tensor and contracting the input
indices.}, label={lst:core-inject}, firstline=76, lastline=97]{code/qrc_core.py}
\end{strip}

The reservoir class (Listing~\ref{lst:core-class}) adds the two feature
amplifiers of Section~\ref{sec:extraction}: temporal multiplexing --- the
interval propagator is the $V$-th root $U_{\Delta t/V}$, applied $V$ times
with $\ev{Z_i}$ read after each sub-step --- and second-moment observables
$\ev{Z_iZ_j}$ read once per interval. With $n=5$ and $V=4$ this yields
$nV + \binom{n}{2} = 30$ features per time step, the count at which all
size-matched classical baselines in this document are set.

\begin{strip}
\lstinputlisting[caption={The reservoir driver in \code{qrc\_core.py}: exact
propagation with $V$ virtual nodes and optional two-point correlators.},
label={lst:core-class}, firstline=102, lastline=142]{code/qrc_core.py}
\end{strip}

\subsection{The validation suite}
\label{sec:validation-suite}

\code{qrc\_core.py} ends with \code{run\_validation\_suite()}, executed before
any experiment in this document was trusted. The suite checks anchors whose
expected values are known independently of the code, so that a silent bug
cannot masquerade as a physics result: (i)~\emph{injection rebuild} --- applying
\code{inject} and tracing back out returns the injected state and the retained
marginal exactly; (ii)~\emph{channel sanity} --- the propagator is unitary to
machine precision and the evolved $\rho$ stays Hermitian, unit-trace, and
positive; (iii)~\emph{ridge recovery} --- on synthetic data generated by a known
linear map with negligible noise, Eq.~\eqref{eq:ridge-solution} recovers the
map to $10^{-10}$; (iv)~\emph{NMSE anchors} --- the mean predictor scores
$1.000$ and the exact-target ``predictor'' scores $0.000$;
(v)~\emph{shot-noise scaling} --- the empirical standard deviation of the
sampled features follows Eq.~\eqref{eq:shot-noise}, halving when $S$
quadruples. Anchors (i)--(v) all pass in the shipped code; the discipline of
\emph{building the anchors before the experiment} is restated as a design
principle in Part~VIII because, in the authors' experience, it is the single
highest-yield habit in quantum machine-learning practice.

\section{Cirq: transparent circuits and exact expectations}
\label{sec:cirq}

Cirq (Google) exposes circuits as plain Python data structures --- moments of
gates on line or grid qubits --- with first-class parameter resolution through
\code{sympy} symbols. That transparency makes it the right teaching vehicle
for the \emph{restart-protocol} reservoir of
Section~\ref{sec:arch-circuit}: one template circuit is built once with
symbolic input angles, and each time step resolves the symbols with the
current input window. Listing~\ref{lst:cirq} is the complete program.

\begin{strip}
\lstinputlisting[caption={\code{cirq\_qrc\_minimal.py}: a complete
restart-protocol quantum reservoir in Cirq. Test NMSE on NARMA-2: $0.277$
(exact expectations, $5$ qubits, window $6$).},
label={lst:cirq}]{code/cirq_qrc_minimal.py}
\end{strip}

\emph{Walk-through.} The configuration block fixes $5$ line qubits, a window
of $L=6$ re-uploaded inputs, and --- decisively, as shown below --- an encoding
scale of $\pi/2$. The template builds, per window slot $j$: one moment of
$R_Y(g_i\,\theta_j)$ input rotations with fixed random per-qubit gains $g_i$
(so different qubits see different projections of the same input), then a
frozen two-layer entangling block of random $R_Y$ rotations and brickwork
$CZ$s --- the reservoir body, drawn once at seed time and never trained.
\code{features\_for\_window} resolves the six symbols against the current
window and calls \code{simulate\_expectation\_values}, which returns
\emph{exact} statevector expectations of the requested observables ---
Cirq's cleanest feature for reservoir research, since it separates
architecture quality from sampling noise by construction. The observable set
is $\ev{Z_i}$, $\ev{X_i}$, and all $\ev{Z_iZ_j}$: twenty features. The
optional \code{SHOTS}$>0$ branch shows the sampled path
(\code{cirq.measure} with key \code{"z"}, then $\ev{Z}=1-2P(1)$ and pairwise
products from the bit table); the $X$ expectations would need a basis change
before measurement, which the exact path sidesteps. Ridge and NMSE close the
file; the washout here discards early steps whose windows are zero-padded.

\emph{The two design lessons the file encodes.} Both were found the honest
way --- by the sweep in Table~\ref{tab:cirq-sweep} --- and generalise well beyond
Cirq. First, \emph{observable richness}: with local $\ev{Z_i}$ alone the
twenty-dimensional dynamics collapse onto five nearly collinear features and
the model barely beats the mean predictor; adding $\ev{X_i}$ and
$\ev{Z_iZ_j}$ triples the usable directions and cuts the error by a factor of
three at the good scale. Second, \emph{encoding scale}: at $\gamma=\pi$ the
re-uploaded rotations drive the features into high harmonics of the input
(Section~\ref{sec:nonlinearity}) that generalise poorly; halving to $\pi/2$
is worth more than any other single change. Neither lesson is visible without
the sweep --- which is the third lesson.

\begin{table}[h]
\centering\small
\caption{Encoding-scale and observable sweep for
Listing~\ref{lst:cirq} (NARMA-2 test NMSE; exact expectations; simulation).
The shipped configuration is bold.}
\label{tab:cirq-sweep}
\begin{tabular}{@{}cccc@{}}
\toprule
Window $L$ & Scale $\gamma$ & Observables (count) & NMSE \\
\midrule
6 & $\pi$   & $Z$ (5)            & 0.948 \\
6 & $\pi$   & $Z,X,ZZ$ (20)      & 0.690 \\
6 & $\pi/2$ & $Z$ (5)            & 0.848 \\
\textbf{6} & $\bm{\pi/2}$ & $\bm{Z,X,ZZ}$ \textbf{(20)} & \textbf{0.277} \\
8 & $\pi$   & $Z,X,ZZ$ (20)      & 0.830 \\
8 & $\pi/2$ & $Z,X,ZZ$ (20)      & 0.402 \\
\bottomrule
\end{tabular}
\end{table}

\section{Qiskit: everything between algorithm and device}
\label{sec:qiskit}

Qiskit (IBM) is the most hardware-aware of the three frameworks, and this
section uses it for exactly that layer: what changes when the reservoir of
Section~\ref{sec:cirq} must run on a machine with a fixed qubit topology, a
native gate alphabet, decoherence, read-out error, and a finite shot budget.
Listing~\ref{lst:qiskit} is the complete program; it is written so that
replacing \code{AerSimulator} by a \code{QiskitRuntimeService} backend --- and
nothing else --- submits the same experiment to an IBM device.

\begin{strip}
\lstinputlisting[caption={\code{qiskit\_qrc\_hw.py}: a hardware-aware windowed
reservoir (Part~A) and a mid-circuit measurement-and-reset demonstration
(Part~B). Simulation with a realistic noise model.},
label={lst:qiskit}]{code/qiskit_qrc_hw.py}
\end{strip}

\subsection{Transpilation: coupling maps and basis gates}
\label{sec:transpilation}

A circuit as written uses idealised gates between arbitrary qubit pairs; a
device executes a fixed native alphabet between physically coupled pairs.
\code{transpile} bridges the two: it \emph{routes} (inserting SWAPs where the
logical circuit touches uncoupled qubits), \emph{translates} into the basis
alphabet --- here IBM's current \code{ecr}, \code{rz}, \code{sx}, \code{x},
\code{id}, where \code{rz} is a software frame change and effectively free ---
and \emph{optimises} at the requested level. The listing pins a linear
five-qubit \code{CouplingMap} to make routing costs visible and
\code{seed\_transpiler} for reproducibility (routing is stochastic). Measured
outcome for one window circuit: depth $78$ with $32$ two-qubit \code{ecr}
gates --- the number that, multiplied by the two-qubit error rate, dominates
the noise budget below. Against a real backend one passes
\code{backend.target} instead of the manual map and lets the reported
calibration data drive layout selection; the manual construction here exists
so the reader can see exactly which constraints bind. The practical rule the
numbers teach: on linear topologies, reservoir circuits should be written
brickwork-local from the start (as this one is), because every nonlocal gate
becomes three.

\subsection{Noise models and finite shots}

The \code{NoiseModel} attaches depolarising errors of $4\times10^{-4}$ per
single-qubit gate and $6\times10^{-3}$ per \code{ecr} --- mid-table values for
current superconducting hardware --- plus an asymmetric read-out confusion
matrix ($1\%$/$2\%$). With $32$ two-qubit gates per circuit the accumulated
depolarisation is already tens of percent of a fully mixed admixture, which
is why the quantum-only column below sits where it does. All $T=400$ step
circuits are transpiled and submitted as \emph{one batched job} of
\code{SHOTS}$=2048$ each --- the shape of a real hardware submission, and the
efficient one, since per-job overhead dwarfs per-circuit cost on queued
devices. \code{counts\_to\_features} converts each counts dictionary to
$\ev{Z_i}$ and $\ev{Z_iZ_j}$; the one genuine trap is bit order ---
\emph{Qiskit keys are little-endian} (qubit $0$ is the rightmost character),
handled here by reversing the key once, in exactly one place.

\subsection{Read-out, honestly regularised, and what the numbers say}
\label{sec:qiskit-results}

The read-out augments the fifteen quantum features with three classical lags
of the input (the hybrid posture of Section~\ref{sec:arch-hybrid}) and
selects the ridge penalty by generalised cross-validation on the training
span --- deterministic and split-free, chosen after a short chronological
validation tail proved unrepresentative on this cloud-regime-switching
series. Table~\ref{tab:qiskit-results} gives the test results on one-step
prediction of the (surrogate) solar clear-sky index.

\begin{table}[h]
\centering\small
\caption{Qiskit pipeline test NMSE (solar clear-sky index, one step ahead;
noisy simulation, $2048$ shots). ``Quantum only'' is the fifteen measured
features; ``classical only'' is three lags of the input.}
\label{tab:qiskit-results}
\begin{tabular}{@{}lc@{}}
\toprule
Feature set & NMSE \\
\midrule
Quantum $+$ classical lags (hybrid) & 0.192 \\
Quantum features only & 0.606 \\
Classical lags only & 0.134 \\
Persistence & 0.138 \\
\bottomrule
\end{tabular}
\end{table}

Read the table the way Part~VIII will insist all such tables be read. The
task is nearly persistent, so persistence and three lags are strong; the
noisy quantum features alone are far worse; and the hybrid \emph{recovers
most but not all} of the classical baseline --- with $\sim$250 training
points, fifteen shot-noisy columns inject enough variance that even a tuned
ridge cannot fully ignore them. There is no quantum benefit here, and the
pipeline says so. The value of the exercise is that every mechanism a
hardware paper must control --- routing overhead, noise accumulation, shot
variance, regularisation under feature noise, and the ease of accidentally
reporting only the flattering column --- is present, measured, and small
enough to fit in one file.

\subsection{Mid-circuit measurement and reset}
\label{sec:midcircuit}

Part~B of the listing demonstrates the dynamic-circuit primitive that
realises Eq.~\eqref{eq:fn-update} on hardware and underlies the NISQRC
result~\citep{hu2024overcoming}: one \emph{long} circuit processes twelve
inputs sequentially; each step encodes $u_k$ on the I/O qubit, entangles,
measures the I/O qubit into \emph{its own single-bit classical register}, and
resets it. Per-step registers matter --- a single register would overwrite the
history --- and the space-joined, reversed key parsing at the end reconstructs
the chronological record. The printed per-step $\ev{Z_{\mathrm{io}}}$
trajectory responds to the input sequence with visible lag: the persistent
qubits are carrying information forward across measurements, on a circuit
whose total duration would exceed any coherence budget if the memory lived in
undisturbed phases. On real IBM backends the same construction runs under
dynamic-circuit support with \code{if\_else} available for the feedback
architectures of Section~\ref{sec:arch-measurement}.

\section{PennyLane: batching and the differentiable boundary}
\label{sec:pennylane}

PennyLane (Xanadu) organises computation around the \emph{QNode} --- a quantum
function bound to a device that behaves as a differentiable array operation
under autograd, JAX, PyTorch, or TensorFlow interfaces. Two capabilities pay
rent in reservoir work: \emph{parameter broadcasting}, which evaluates a
whole time series of input windows in one batched call, and
\emph{differentiability}, which is deliberately used here for one scalar ---
the encoding gain --- while the reservoir body stays frozen.
Listing~\ref{lst:pennylane} is the complete program; the task is 3-month-ahead
prediction of the \emph{real} ENSO sea-surface-temperature anomaly.

\begin{strip}
\lstinputlisting[caption={\code{pennylane\_qrc\_hybrid.py}: a broadcast
windowed quantum feature map with a trained ridge read-out, plus alternating
gradient tuning of a single encoding gain. Real ENSO data.},
label={lst:pennylane}]{code/pennylane_qrc_hybrid.py}
\end{strip}

\emph{Walk-through.} The QNode re-uploads a window of $8$ inputs through
gain-weighted $R_Y$ rotations interleaved with a frozen
\code{StronglyEntanglingLayers} block, and returns twenty expectations
($Z$, $X$, $ZZ$). The ellipsis indexing \code{u\_win[..., j]} is the
broadcasting idiom: called with a $(T,8)$ array, the device vectorises over
the leading axis and one call returns all $T$ feature rows. The tuning stage
then alternates two steps: refit the ridge read-out classically at the
current gain $\gamma$; take one gradient step on $\gamma$ with that read-out
\emph{frozen}, so the gradient flows only through the quantum circuit and no
differentiation through the linear solve is attempted --- the standard, stable
recipe at the reservoir/variational boundary, and vastly cheaper than
end-to-end training since the quantum parameter count is one.

\emph{What the numbers teach.} Fixed $\gamma=1.0$ scores NMSE $0.976$ ---
barely better than the mean; fourteen alternating iterations walk $\gamma$
to $0.69$ and the error to $0.747$; persistence scores $0.695$; and the
\emph{recurrent} exact reservoir of Section~\ref{sec:reference-impl} on the
same task scores $0.596$ (Part~VII). The ranking is the lesson. A windowed
feature map, however tuned, hard-truncates memory at $L=8$ months, and ENSO
carries predictive structure well beyond that; genuine recurrence --- state
carried forward indefinitely with fading weights --- is not a nicety but the
difference between beating persistence and not. Differentiability recovered
what the encoding scale had thrown away (a threefold error reduction, for
one parameter and fourteen circuit-pair evaluations) but could not
manufacture memory the architecture does not have. When to cross fully into
variational training --- trainable body, parameter-shift gradients through
hundreds of parameters --- is a cost question: each shift-rule gradient costs
two full batched evaluations \emph{per parameter}, which is precisely the
expense the reservoir paradigm exists to avoid
\citep{mitarai2018quantum,cerezo2021variational}.

\section{Shot noise in practice}
\label{sec:shot-practice}

The exact simulations above have an analytic ceiling; hardware estimates
approach it as $S^{-1/2}$. The shipped shot-noise model perturbs each exact
Pauli expectation $z$ by binomially exact noise,
$\hat z = \mathrm{clip}\bigl(z + \xi\sqrt{(1-z^2)/S},\,-1,\,1\bigr)$ with
$\xi\sim\mathcal N(0,1)$, and the validation suite confirms the empirical
scaling of Eq.~\eqref{eq:shot-noise}. Applied to the NARMA-10 reference task
(Part~VII, Figure~\ref{fig:shot-noise}): the exact-feature ceiling is NMSE
$0.283$; at $S=10^2$ the score is $0.90$, at $S=10^3$ it is $0.64$, at
$S=10^4$ it is $0.53$, and even $S=10^5$ leaves $0.45$ --- still $60\%$ above
ceiling. Three practices follow, used throughout this document and restated
in Part~VIII: measure the ceiling first, so architecture and sampling effects
are never conflated; report $S$ next to every number, since shots are the
real hardware currency; and let the regulariser know --- feature noise of known
variance shifts the optimal $\lambda$ upward, which is why every read-out in
this document selects $\lambda$ on data rather than by convention. The
fundamental version of these observations --- sampling noise caps the number
of resolvable feature directions (\emph{eigentasks}), with an
information-optimal read-out construction --- is developed in
\citet{hu2023tackling}, and shot-optimal training specifically in
\citet{ahmed2025optimal}.

\section{Choosing a framework}
\label{sec:framework-choice}

Table~\ref{tab:framework-compare} condenses the working differences; the
recommendation logic beneath it is short. For \emph{architecture research} ---
new encodings, observables, dynamical regimes --- use the exact NumPy path (or
Cirq's exact expectations) so sampling never confounds design, and port later.
For \emph{hardware-bound work} on superconducting devices, build in Qiskit
from day one: transpilation realism, calibration-driven noise models, and
dynamic circuits are where that project's risk lives. For \emph{hybrid
models} that may drift toward trainable components, or that live inside a
JAX/PyTorch stack, PennyLane's interfaces and broadcasting make it the
lowest-friction host. The three are not exclusive: this document's own
pipeline prototyped exactly, validated in Cirq, and hardened in Qiskit.

\begin{table*}[tp]
\centering\small
\caption{Working comparison of the three primary frameworks for QRC
(versions as used here: Cirq 1.7, Qiskit 2.5/Aer 0.17, PennyLane 0.45).}
\label{tab:framework-compare}
\begin{tabular}{@{}p{3.4cm}p{3.6cm}p{3.6cm}p{3.6cm}@{}}
\toprule
Axis & Cirq & Qiskit & PennyLane \\
\midrule
Circuit model & moments of gates; explicit, list-like & \code{QuantumCircuit}; register-oriented & quantum functions (QNodes) \\
Exact expectations & \code{simulate\_expectation}\newline\code{\_values} (statevector) & \code{Estimator} primitives; Aer & analytic on \code{default.qubit} \\
Density-matrix / noise & \code{DensityMatrixSimulator}; per-moment noise & Aer \code{NoiseModel}: gate errors, read-out, thermal; calibration import & \code{default.mixed}; channel ops \\
Transpilation to real topologies & basic device gatesets & \textbf{best in class}: routing, basis translation, optimisation levels, targets & delegates to plugins \\
Mid-circuit measure/reset & supported in simulators & \textbf{production dynamic circuits}; per-step classical registers (Listing~\ref{lst:qiskit}) & \code{qml.measure(reset=True)}; allocates an extra wire on \code{default.mixed} --- prefer compiled-unitary or NumPy paths for density-matrix studies \\
Batching a time series & loop over resolvers (fast) & batched job of many circuits (hardware-shaped) & \textbf{parameter broadcasting}: one call \\
Differentiability & external & external & \textbf{native}: autograd/JAX/Torch/TF, parameter shift \\
Symbolic parameters & \code{sympy}, first-class & \code{Parameter} objects & function arguments \\
Hardware access & Google-adjacent; via converters & IBM Quantum; broad vendor reach & vendor plugins (IBM, IonQ, Braket, \ldots) \\
Best QRC role here & transparent restart-protocol teaching; exact studies & the hardware-facing pipeline & hybrid/differentiable boundary; fast batched sweeps \\
\bottomrule
\end{tabular}
\end{table*}

\begin{table*}[h]
\centering\small
\caption{Secondary tools, one line each, for completeness.}
\label{tab:secondary-tools}
\begin{tabular}{@{}lp{10.6cm}@{}}
\toprule
Tool & One-line role for QRC work \\
\midrule
QuTiP~\citep{johansson2012qutip} & the open-systems workhorse: Lindblad solvers for dissipative reservoirs (Section~\ref{sec:arch-dissipative}) \\
Qulacs~\citep{suzuki2021qulacs} & fastest CPU statevector loops for large sweep campaigns \\
TensorCircuit~\citep{zhang2023tensorcircuit} & tensor-network backends and JAX-native batching for larger $n$ at low entanglement \\
CUDA-Q / cuQuantum & GPU statevector and density-matrix kernels when sweeps outgrow CPUs \\
Strawberry Fields~\citep{killoran2019strawberry} & the CV/photonic reservoirs of Section~\ref{sec:arch-cv}, Gaussian and Fock backends \\
Amazon Braket & multi-vendor hardware access (IonQ, Rigetti, QuEra) behind one API \\
Quantinuum toolchain & H-series access; mid-circuit measurement and QCCD primitives \\
IonQ / Rigetti clouds & direct vendor access for the platforms of Part~V \\
Mitiq~\citep{larose2022mitiq} & framework-agnostic error mitigation (ZNE, PEC) for feature estimation \\
\bottomrule
\end{tabular}
\end{table*}
